\subsection{Operações de proximidade em estações espaciais}
\label{subsec_operacoes_proximidade_estacoes}

A evolução das operações de \gls{rvdb} possibilitou a realização de missões logísticas e de manutenção em estações espaciais, nas quais veículos visitantes realizam aproximações controladas para transporte de tripulação, suprimentos e equipamentos científicos. 
Essas operações exigem elevada precisão na navegação relativa e no controle de atitude durante as fases finais de aproximação, uma vez que pequenas incertezas na estimativa de posição ou orientação podem comprometer a segurança do encontro orbital.

A Estação Espacial Internacional (\gls{iss}) representa um dos principais exemplos de infraestrutura orbital que depende de operações recorrentes de proximidade. 
Desde o início de sua montagem, múltiplos veículos visitantes passaram a realizar missões de reabastecimento e transporte de tripulação, exigindo o desenvolvimento de diferentes estratégias de aproximação e acoplamento.

Além das operações de acoplamento direto (\emph{docking}), algumas missões utilizam o processo de atracação assistida (\emph{berthing}), no qual o veículo visitante é capturado por um \gls{mr} e posteriormente conduzido para acoplar-se à estação espacial. 
Nesse tipo de operação, o veículo aproxima-se da estação até uma posição de captura segura, permanecendo estacionário enquanto o manipulador realiza a captura e conduz o veículo até a interface de conexão.

Entre os exemplos de veículos que utilizam esse tipo de operação está a espaçonave cargueira \emph{Cygnus}, desenvolvida pela empresa \emph{Northrop Grumman} para transportar suprimentos e experimentos científicos para a \gls{iss}. 
Durante suas missões, a \emph{Cygnus} realiza a aproximação até uma posição de captura, sendo então agarrada pelo \gls{mr} \emph{Canadarm2}, operado pelos astronautas a bordo da estação, que finalizam o processo de atracação conectando o veículo a um dos módulos da estação \cite{nasaCygnus}.

De forma semelhante, o veículo japonês \emph{H-II Transfer Vehicle} (HTV), desenvolvido pela \gls{jaxa}, também foi projetado para realizar operações de \emph{berthing}. 
Após realizar a aproximação controlada à estação espacial, o HTV era capturado pelo \gls{mr} e posteriormente conectado à estrutura da \gls{iss}, permitindo o acesso às cargas transportadas \cite{nasaHTV}.

O \gls{mr} \emph{Canadarm2}, desenvolvido pela \gls{csa}, desempenha papel central nessas operações. 
Com múltiplos graus de liberdade e capacidade de ser operado tanto por astronautas quanto por controladores em solo, esse sistema permite capturar veículos visitantes, reposicionar módulos da estação e realizar tarefas de manutenção externa \cite{canadarm2}.
